\subsection{中微子--电子散射与反冲谱}
\label{sec:sm-neutrino-electron-applications}

$\mu$ 子寿命固定了低能带电流的强度，中微子--电子弹性散射进一步分辨电子的左右手弱耦合。靶粒子是轻子，因而无需引入强子形状因子。实验却常在电子质量不可忽略的能量区间测量反冲电子，所以应保留精确的二体运动学，再判断能否作高能展开。

考虑 $\nu_a(k)e^-(p)\to\nu_a(k')e^-(p')$，味指标 $a=e,\mu,\tau$。实验室系初始电子静止，入射中微子能量为 $E_\nu$，末态电子动能为 $T_e$。定义 $E_{e0}=m_{\mathrm e} c^2$。中微子质量相对 $E_\nu$ 可忽略；电子质量保持有限。完整树级标准模型振幅含 $Z$ 交换，$a=e$ 时还含 $W$ 交换。只有在所有相关传递能量平方远小于 $M_W^2,M_Z^2$ 时，才把传播子展开成局域作用，其中 $M_V=m_Vc^2$ 是规范玻色子的静止能量。

为书写无量纲约化振幅，定义能量四矢量 $K=ck$、$P=cp$、$K'=ck'$、$P'=cp'$。这只是带显式 $c$ 的变量替换，并未设置 $c=1$；相应旋量满足 $\sum_s u(P,s)\bar u(P,s)=\slashed P+E_{e0}$。四费米最低阶振幅可写为
\begin{equation}
 \begin{aligned}
 \mathcal M_{\nu_a e}
 &=2\sqrt2G_F^{(E)}
 [\bar u(K')\gamma^\mu P_Lu(K)]
 [\bar u(P')\gamma_\mu(g_L^{(a)}P_L+g_R^{(a)}P_R)u(P)],\\
 g_L^{(a)}&=-\frac12+s_W^2+\delta_{ae},\qquad
 g_R^{(a)}=s_W^2.
 \end{aligned}
 \label{eq:sm-nue-fermi-amplitude}
\end{equation}
一个公共整体相位不影响下文截面。$s_W=\sin\theta_W$，$g_L^{(a)},g_R^{(a)}$ 是本过程的无量纲电子手征系数；它们不包含公共 Fermi 常数。$\delta_{ae}$ 是 Kronecker 符号，来自电子味的带电流；$W$ 图经左手流 Fierz 重排后与 $Z$ 图具有相同外腿配对，因此应在振幅中相加。最低阶中性流值也可由 $g_L=(g_V^e+g_A^e)/2$、$g_R=(g_V^e-g_A^e)/2$ 检查。精密分析中这些树级系数需要按指定重整化方案加入辐射修正\cite{ErlerFreitas2024EW}。

末态中微子在壳条件给出电子反冲的精确范围
\begin{equation}
 0\leq T_e\leq T_{\max},\qquad
 T_{\max}=\frac{2E_\nu^2}{E_{e0}+2E_\nu}.
 \label{eq:sm-nue-recoil-endpoint}
\end{equation}
在上述四费米近似内，非极化电子靶的微分截面为
\begin{equation}
 \frac{\dd\sigma_{\nu_a e}}{\dd T_e}
 =\frac{2[G_F^{(E)}]^2(\hbar c)^2E_{e0}}{\pi}
 \left[(g_L^{(a)})^2+(g_R^{(a)})^2\left(1-\frac{T_e}{E_\nu}\right)^2
 -g_L^{(a)}g_R^{(a)}\frac{E_{e0}T_e}{E_\nu^2}\right].
 \label{eq:sm-nue-recoil-spectrum}
\end{equation}
反中微子散射由方括号内 $g_L^{(a)}\leftrightarrow g_R^{(a)}$ 得到。只有初始电子的两个自旋需要平均；入射中微子束只有参与弱相互作用的螺旋度，不能再除以二。前因子的单位为 $\mathrm{m^2/J}$，因此反冲能量积分恢复面积。末项是有限电子质量引起的左右手干涉，在低能区与另外两项同样重要。

令 $C_{\nu e}=2[G_F^{(E)}]^2(\hbar c)^2E_{e0}/\pi$，对任意真实系数 $g_L,g_R$ 定义
\begin{equation}
 \mathcal I(T;g_L,g_R)
 =C_{\nu e}\left[g_L^2T+g_R^2\left(T-\frac{T^2}{E_\nu}
 +\frac{T^3}{3E_\nu^2}\right)-g_Lg_R\frac{E_{e0}T^2}{2E_\nu^2}\right].
 \label{eq:sm-nue-integrated-spectrum}
\end{equation}
于是 $\sigma=\mathcal I(T_{\max};g_L,g_R)$；具有理想反冲阈值 $T_{\min}<T_{\max}$ 的截面为 $\mathcal I(T_{\max})-\mathcal I(T_{\min})$。当阈值高于运动学端点时，截面为零。真实探测效率依赖 $T_e$ 时，须在积分内乘 $\mathcal Q^{\rm det}(T_e)$。

只有在 $E_\nu\gg E_{e0}$ 时，才可以令 $T_{\max}\simeq E_\nu$ 并舍弃相对 $E_{e0}/E_\nu$ 修正，得到
\begin{equation}
 \sigma_{\nu_a e}\simeq C_{\nu e}E_\nu
 \left[(g_L^{(a)})^2+\frac13(g_R^{(a)})^2\right].
 \label{eq:sm-nue-high-energy}
\end{equation}
低能极限 $E_\nu\ll E_{e0}$ 则有 $T_{\max}\simeq2E_\nu^2/E_{e0}$，截面随 $E_\nu^2$ 而非 $E_\nu$ 增长。此时积分中的质量干涉项不能省略：最低阶组合为 $g_L^2+g_R^2-g_Lg_R$，它在交换中微子与反中微子后不变。

\begin{exbox}[有限质量与阈值的影响]
取教学参数 $E_\nu=\SI{1}{\mega\electronvolt}$、$E_{e0}=\SI{0.510999}{\mega\electronvolt}$、$s_W^2=0.231$，以及 $G_F^{(E)}=1.16638\times10^{-5}\,\mathrm{GeV^{-2}}$。求电子味中微子的反冲端点、总截面及给定阈值以上的接收截面，并与高能极限比较。
\textbf{解。} 电子味中微子有 $g_L=0.731,g_R=0.231$，反冲端点为
\begin{equation*}
 T_{\max}=\SI{0.79650}{\mega\electronvolt}.
\end{equation*}
将所有能量统一以 GeV 代入\eqnrefs{eq:sm-nue-integrated-spectrum}，再用 $(\hbar c)^2$ 换成面积，得到总截面约 $7.17\times10^{-49}\,\mathrm{m^2}$；高能公式给 $9.52\times10^{-49}\,\mathrm{m^2}$，高估约 $33\%$。若只接收 $T_e>\SI{0.2}{\mega\electronvolt}$，可见截面约为 $5.21\times10^{-49}\,\mathrm{m^2}$，接收比例约 $0.726$。

上述 $s_W^2$ 是用来演示计算的固定树级输入，不是对 MeV 实验有效弱混合角的精密拟合。原子电子可近似为自由静止电子还要求反冲显著高于相关束缚能；接近原子激发或电离阈值时，需要原子响应函数。
\end{exbox}

\begin{figure}[htbp]
 \centering
 \includegraphics[width=0.75\textwidth]{neutrino_electron_recoil.pdf}
 \caption{给定 $E_\nu=\SI{1}{\mega\electronvolt}$ 时电子味和 $\mu$ 味中微子、反中微子的电子反冲谱。纵轴除以共同前因子 $C_{\nu e}$；所有曲线在同一有限质量端点终止。}
 \label{fig:sm-nue-recoil}
\end{figure}
\figref{fig:sm-nue-recoil} 中电子味的额外带电流改变谱的归一化和形状；中微子与反中微子的比较则交换两种手征权重。这是同一局域弱作用在不同束流条件下的可测区别。

